% Spatial Ecology in R Presentation in LaTeX
\documentclass[8pt]{beamer}

\usepackage{helvet}
\renewcommand{\familydefault}{\sfdefault}


% Theme choice:
\usetheme{Frankfurt}
%
%
%change colour of theme
\usecolortheme{seahorse}
%
% Packages
\usepackage[utf8]{inputenc}  % For UTF-8 encoding
\usepackage{graphicx}        % For including images
\usepackage{amsmath}         % For mathematical symbols and equations



\title{Spatial Ecology in R Final Examination Project}
\subtitle{Vegetation Analysis in Northwestern Thailand Using R}
\institute{University of Bologna}
\author{Sofiya Tumanova}
\date{\today}

\begin{document}


\begin{frame}
    \titlepage
\end{frame}

% Table of Contents
\begin{frame}{Table of Contents}
    \tableofcontents
\end{frame}

% Section 1
\section{Introduction}
\subsection{Project Overview}
\begin{frame}{Project Overview}
    \textbf{Overview}\\
    Multi temporal analysis of vegetation in Northwestern Thailand. Specifically focusing on Fraction of Green Vegetation Cover and the Normalized Difference Vegetation Index.\\ 
    \smallskip
    \textbf{Objectives}
    \begin{itemize}
        \item Analyzing FCOVER changes over a 20 year period\\
        (1999 - 2019)
        \item Performing a classification of the FCOVER
        \item Computing the vegetation index (NDVI)
        \item Classifying the NDVI raster output
        \item Analyzing the differences in NDVI over a 7 year period\\
        (2016 - 2023)
    \end{itemize}    
\end{frame}

\begin{frame}{Vegetation Analysis}
    \begin{columns}[T] % Align columns at the top
        \column{0.5\textwidth}
        \textbf{FCOVER}
        \begin{itemize}
            \item Fraction of Green Vegetation Cover
            \item Represents the portion of ground covered by green vegetation
            \item \textbf{Range:} 0 - 1
                \begin{itemize}
                    \item 0: No Vegetation Cover (Bare Soil, Water)
                    \item 1: Full Vegetation Cover
                \end{itemize}
            \item \textbf{Common Uses:}
                \begin{itemize}
                    \item Land Use and Land Cover
                    \item Forestry - estimate canopy density
                    \item Climate Studies
                \end{itemize}
        \end{itemize}

        \column{0.5\textwidth}
        \textbf{NDVI}
        \begin{itemize}
            \item Normalized Difference Vegetation Index
            \item Measures vegetation health and density
            \item Uses reflectance values from the RED and NIR bands
            \item \textbf{Range:} -1 to 1
                \begin{itemize}
                    \item -1: Water or No Vegetation
                    \item 1: Dense Vegetation
                \end{itemize}
            \item \textbf{Common Uses:}
                \begin{itemize}
                    \item Vegetation Health
                    \item Agriculture - crop stress
                    \item Drought Monitoring
                \end{itemize}
        \end{itemize}
        \begin{figure}
            \centering
            \includegraphics[width=0.5\textwidth]{NDVI stressed vs healthy leaves.jpg}
        \end{figure}
    \end{columns}
\end{frame}

\subsection{Areas of Interest}
\begin{frame}{Areas of Interest} 
    \begin{columns}[T]
        \column{0.5\textwidth}
        \begin{itemize}
            \item \textbf{Study Area 1:}\\ FCOVER Analysis of Northwestern Thailand: 1999 to 2019\\
            98.5, 100.7, 16.6, 18.9 ( 59,394 km\(^2\))
            \item \textbf{Study Area 2:}\\ NDVI Analysis of Doi Inthanon National Park and Surrounding Areas: 2016 and 2023\\
             98.3, 99.4, 17.9, 19.0 (13,788 km\(^2\)) 
        \end{itemize}

        \column{0.5\textwidth}
        \centering
        \includegraphics[width=0.9\textwidth]{Spatial ecology in R.png}
    \end{columns}
\end{frame} 
   


% Section 2
\section{Data}
\subsection{Data Download}
\begin{frame}{Data Download} 
    \begin{columns}[t] % Aligns the top of the columns
        \column{0.5\textwidth} 
        \textbf{FCOVER}
        \begin{itemize}
            \item Copernicus Land Monitoring Service\\
            February: 1999, 2004, 2009, 2014, 2019
            \item Spatial Resolution: 1km
            \item Temporal Resolution: 10 days
            \item \url{https://land.copernicus.eu/en}
        \end{itemize}

        \column{0.5\textwidth} 
        \textbf{NDVI}
        \begin{itemize}
            \item Copernicus Browser\\
            November: 2016, 2023
            \item Satellite: Sentinel 2
            \item Level: 2A BOA
            \item Cloud Cover less than 10\%
            \item Bands: 2, 3, 4, 8
            \item \url{https://browser.dataspace.copernicus.eu/}
        \end{itemize}
    \end{columns}
\end{frame}

\subsection{Pre-Processing}
\begin{frame}{Pre-Processing: Fraction of Green Vegetation Cover}
\begin{columns}[t]
        \column{0.3\textwidth} 
        \textbf{FCOVER}
        \begin{itemize}
            \item Data was loaded into R and cropped to the same extent using the \texttt{crop} function from the terra package\\
            \item Data was visualized by plotting each raster with \texttt{ggplot}
            \item Viridis option "mako" was used
        \end{itemize}

        \column{0.7\textwidth} 
        \textbf{FCOVER Plots for a 20 Year Period}
        \centering
        \includegraphics[width=0.9\textwidth]{FCOVER 1999 to 2019 plot reclassified.jpeg}
        
    \end{columns}
\end{frame}

\begin{frame}{Pre-Processing: Normalized Difference Vegetation Index}
\small
\begin{columns}[t]
        \column{0.3\textwidth} 
        \textbf{NDVI}
        \begin{itemize}
            \item Each band was imported into R Studio and cropped to the same extent using the \texttt{crop} function\\
            \item A true color and false color composite was created for visualization purposes
            \item True Color Composite: R, G, B (4,3,2)
            \item False Color Composite: NIR, R, G (8,4,3)
        \end{itemize}

        \column{0.7\textwidth} 
    \textbf{True and False Color Composites}
    \centering
    \vbox{
        \includegraphics[width=0.9\textwidth]{color composite 2016.png}\\
        \includegraphics[width=0.9\textwidth]{color composite 2023.png}
    }
    \end{columns}
\end{frame}



% Section 3
\section{Analysis and Results}
\subsection{FCOVER Analysis}

\begin{frame}{FCOVER Analysis: Difference Raster}
\small  % Reduce font size for the entire slide
\begin{columns}[T]
    \column{0.3\textwidth} 
    \textbf{20-year difference in FCOVER}
    \begin{itemize}
        \item To visualize change, the 1999 FCOVER raster was subtracted from the 2019 FCOVER raster.
        \item The output raster was plotted with \texttt{ggplot} function
        \item Viridis color scale option "magma" was used for visualization.
    \end{itemize}

    \column{0.7\textwidth} 
    \centering
    \includegraphics[width=0.9\textwidth]{FCOVER diff from 1999 to 2019.png}
    
\end{columns}
\end{frame}


\begin{frame}{FCOVER Analysis: Classification}
  \small  % Reduce font size for the entire slide
\begin{columns}[T]
    \column{0.4\textwidth} 
    \textbf{Classification of FCOVER difference raster}
    \begin{itemize}
        \item For better interpretation, the difference raster was classified into 3 categories\\
        Decrease, Increase, No Change
        \item The \texttt{im.classify} function from the imageRy package was used for classification.
        \item Threshold for classes : -0.0001,0,0.0001
    \end{itemize}

    \column{0.6\textwidth} 
    \centering
    \includegraphics[width=0.9\textwidth]{FCOVER Change Classified.jpeg}
    
\end{columns} 
\end{frame}

\begin{frame}{FCOVER Analysis: Side by Side Comparison}
\centering
\includegraphics[width=0.9\textwidth]{FCOVER change and classified side by side.png}
    
\end{frame}

\begin{frame}{FCOVER Analysis: Calculation of Percent Change}
\textbf{Methods:}
    \begin{itemize}
        \item Counting total number of pixels using \texttt{ncell} function from the terra package
        \item Counting the percentage of pixels falling into each category with the following formula:\\
        \[
\frac{\text{freq}(\text{class.diff.1999.2019})}{\text{pixel\_total}} \times 100
\]
        \item Creating a vector with 3 categories:\\
         Decrease, No Change, Increase
        \item Creating a second vector with the percentages calculated earlier
        \item Creating a dataframe and populating the values by creating two vectors, one with the 3 change categories, and another with the calculated percent values
        \item Plotting using \texttt{ggplot}
    \end{itemize}  
\end{frame}

\subsection{FCOVER Results}

\begin{frame}{Results FCOVER Analysis: Calculation of Percent Change}
    %\centering
    \includegraphics[width=0.9\textwidth]{FCOVER change in percentages.jpeg}
\end{frame}


\subsection{NDVI Analysis}

\begin{frame}{NDVI Analysis: Calculating NDVI}
\begin{columns}[t]
    \column{0.4\textwidth}
    \textbf{Methods}
    \begin{itemize}
        \item The \textbf{Normalized Difference Vegetation Index} (NDVI) was calculated using the following formula:
        \begin{equation}
            \text{NDVI} = \frac{\text{NIR} - \text{RED}}{\text{NIR} + \text{RED}}
        \end{equation}
        \item Where NIR is the Near Infrared Band (Band 8) and RED is the Red Band (Band 4).
        \item This process was repeated with the 2016 and 2023 images.
        \item The NDVI results were plotted using the \texttt{ggplot} function using the "viridis" color palette.
    \end{itemize} 
    
    \column{0.6\textwidth}
    \textbf{NDVI Plots for 2016 and 2023}
    \centering
    \includegraphics[width=0.7\textwidth]{NDVI cropped plots.png}
\end{columns}
\end{frame}

\begin{frame}{NDVI Analysis: Difference Raster}
\small  % Reduce font size for the entire slide
\begin{columns}[T]
    \column{0.4\textwidth} 
    \textbf{7-year difference in NDVI}
    \begin{itemize}
        \item To visualize change, the 2016 NDVI raster was subtracted from the 2023 NDVI raster.
        \item The output raster was plotted with \texttt{ggplot}.
        \item Viridis color scale option "viridis" was used for visualization.
    \end{itemize}

    \column{0.6\textwidth} 
    \textbf{NDVI Difference Plot}
    %\centering
    \includegraphics[width=0.9\textwidth]{NDVI difference 2016 to 2023.jpeg}
    
\end{columns}
\end{frame}

\begin{frame}{NDVI Analysis: Classification of 2016 and 2023 rasters}
\textbf{Methods:}
    \begin{itemize}
        \item Classification of NDVI was based on the study: \\
        \bigskip
        \textit{Urban Vegetation Classification with NDVI Threshold Value Method Using Very High Resolution (VHR) Pleiades Imagery} by Hashim et al., 2019 \\
        \bigskip
        \url{https://isprs-archives.copernicus.org/articles/XLII-4-W16/237/2019/}
        \smallskip
        \item Created a reclassification matrix using the \texttt{matrix} function
        \item Reclassified each NDVI raster using the \texttt{reclassify} function according to the categories:
        \begin{itemize}
            \item \textbf{Non Vegetation} for level 1: \hspace{0.5cm} \(-1.0 \leq \text{NDVI} < 0.2\)
            \item \textbf{Low Vegetation} for level 2: \hspace{0.5cm} \(0.2 \leq \text{NDVI} \leq 0.5\)
            \item \textbf{Dense Vegetation} for level 3: \hspace{0.5cm} \(0.501 \leq \text{NDVI} \leq 1.0\)
        \end{itemize}
        \item Generated plots using the \texttt{ggplot} function
    \end{itemize}  
\end{frame}

\begin{frame}{Classified NDVI Plots 2016 and 2023}
    \centering
    \includegraphics[width=\textwidth, height=\textheight, keepaspectratio]{classified NDVI plots side by side.png}
\end{frame}


\begin{frame}{NDVI Analysis: Difference in Pixel Count Between 2016 and 2023}
\textbf{Methods:}
    \begin{itemize}
        \item Counted the pixels belonging into each NDVI\_Class for 2016 and 2023 classified NDVI rasters 
        \item Functions used: \texttt{n()} and \texttt{summarize()}
        \item Joined counts together into a single df using the \texttt{full\_join} function
        \item Created a “Difference” column and calculated the difference by subtracting 2016 from the 2023 pixel count
        \item Converted the NDVI\_Class column into a factor (categorical variable) using the \texttt{factor} function, which allows me to represent data in a specific order (levels: 1, 2, 3) “Non vegetation”, “Low Vegetation”, and “Dense Vegetation”
        \item Generated plot using the \texttt{ggplot} function
    \end{itemize}
\end{frame}

\subsection{NDVI Results}

\begin{frame}{Results NDVI Analysis: Difference in Pixel Count Between 2016 and 2023}
        \centering
        \includegraphics[width=\textwidth]{difference in pixel counts from 2016 to 2023.png}   
\end{frame}

\begin{frame}{NDVI Analysis: Pixel Count Percent Change 2016 and 2023}
\textbf{Methods:}
    \begin{itemize}
        \item A new df "counts\_summary" was created by combining pixel counts from each df using the function \texttt{merge} based on the common column: "NDVI\_Class"
        \item Percentage change was calculated with the following formula\\
        \[
\text{Percent Change} = \left(\frac{\text{Count}_{2023} - \text{Count}_{2016}}{\text{Count}_{2016}}\right) \times 100
\]
        \item A new column "Percent\_Change" was created with the values obtained from the formula
        \item Generated plot using the \texttt{ggplot} function
    \end{itemize}
\end{frame}

\begin{frame}{Result NDVI Analysis: Pixel Count Percent Change 2016 and 2023}
        \centering
        \includegraphics[width=\textwidth]{percentage change in NDVI categories.png} 
\end{frame}


% Section 4
\section{Conclusion}


\begin{frame}{Summary of Results and Considerations}
\begin{columns}[T]
    \column{0.5\textwidth}
    \begin{itemize}
    \item \textbf{Main Findings FCOVER between 1999 and 2019}
        \begin{itemize}
         \item 13.7\% of pixels had a decrease in FCOVER
         \item 42.21\% of pixels had an increase in FCOVER
         \item 44.09\% of pixels had no change in FCOVER
    \end{itemize}
    \item \textbf{Main Findings NDVI between 2016 and 2023}
    \begin{itemize}
        \item 44.22\% increase of pixels in the "Non Vegetation Category"
        \item 19.62\% increase of pixels in the "Non Vegetation Category"
        \item 2.12\% decrease of pixels in the "Dense Vegetation Category"
    \end{itemize}
    \end{itemize}
        \column{0.5\textwidth}
        \begin{itemize}
            \item \textbf{Suggestions and Future Considerations}
            \begin{itemize}
                \item Removing NA values can have an effect on the overall statistics
                \item Cloud cover influences results
                \item Bigger temporal extent needed for more analysis
                \item Combination with other analysis such as:
                \begin{itemize}
                    \item Soil Moisture Index (SMI)
                    \item Normalized Difference Water Index (NDWI)
                \end{itemize}
            \end{itemize}
        \end{itemize}
    \end{columns}
\end{frame}

% Final frame
\begin{frame}
    \centering
    \Huge{Thank You for Your Attention}\\
    \vspace*{\fill}
    \small
    Universita Di Bologna\\
    sofiya.tumanova@studio.unibo.it
\end{frame}

\end{document}
